\documentclass[11pt]{article}
\usepackage[margin=1in]{geometry}
\usepackage{amsmath,amssymb,amsthm}
\usepackage{graphicx}
\usepackage{hyperref}

\newtheorem{theorem}{Theorem}
\newtheorem{lemma}{Lemma}
\newtheorem{remark}{Remark}

\newcommand{\Tr}{\operatorname{Tr}}
\newcommand{\E}{\mathbb E}
\newcommand{\eps}{\varepsilon}
\newcommand{\C}{\mathbb C}
\newcommand{\diag}{\operatorname{diag}}

\title{A partial sharp Rademacher case for the \\ noncommutative Khintchine constant}
\author{Run fa\_banach\_001, lane 8}
\date{2026-07-18}

\begin{document}
\maketitle

\section*{Status}

This is a partial result for the open sharp-constant question in
Haagerup--Musat's paper \cite{HM}.  The general Rademacher case remains
open.  The result below proves that the conjectural lower constant
\(1/\sqrt2\) is valid for all commuting normal coefficient families and
for arbitrary two-term matrix coefficient families.

\section*{Source evidence}

Haagerup--Musat prove the Rademacher version with lower constant
\(1/\sqrt3\), and state that they do not know whether this is sharp.
The relevant passage is on page 2 of \cite{HM}.

\begin{center}
\includegraphics[width=\linewidth]{figures/open_problem_crop.png}
\end{center}

\section*{Source problem}

For \(x_1,\ldots,x_d\in M_n(\C)\), set
\[
N(x_1,\ldots,x_d)=
\inf_{x_i=y_i+z_i}
\Tr\left(\sum_i y_i^*y_i\right)^{1/2}
+\Tr\left(\sum_i z_i z_i^*\right)^{1/2}.
\]
Haagerup--Musat prove
\[
\frac1{\sqrt3} N(x_1,\ldots,x_d)
\leq
\E_\eps\left\|\sum_i \eps_i x_i\right\|_{S_1^n}
\leq
N(x_1,\ldots,x_d),
\]
where \((\eps_i)\) are independent Rademacher signs.  They also record
the scalar obstruction showing that the best possible lower constant is
at most \(1/\sqrt2\).  The open sharpness issue is whether the lower
constant can be improved from \(1/\sqrt3\) to \(1/\sqrt2\) in full
generality.

\section*{Proof intuition}

The two subcases below remove different sources of noncommutativity.
For commuting normal matrices, simultaneous diagonalization reduces both
the trace norm and the \(N\)-norm to pointwise scalar quantities; Szarek's
sharp scalar \(p=1\) Khintchine inequality then applies on each diagonal
fiber.  For two arbitrary matrices, a \(2\times2\) Hadamard rotation of
the row block \([x_1\ x_2]\) gives the needed comparison with
\(\|x_1+x_2\|_1+\|x_1-x_2\|_1\), which is exactly twice the Rademacher
average.  Thus the remaining obstruction, if any, must use at least
three noncommuting coefficients.

\begin{lemma}[Diagonal formula]
\label{lem:diag}
Suppose \(x_i=\diag(a_i(1),\ldots,a_i(n))\) are diagonal matrices.  Then
\[
N(x_1,\ldots,x_d)
=
\sum_{k=1}^n\left(\sum_{i=1}^d |a_i(k)|^2\right)^{1/2}.
\]
\end{lemma}

\begin{proof}
The upper bound follows by taking \(y_i=x_i\), \(z_i=0\).  For the lower
bound, let \(D:M_n(\C)\to M_n(\C)\) be the diagonal conditional
expectation.  It is the average of conjugations by diagonal sign
unitaries.  Hence the maps
\[
(u_i)_i\mapsto (D u_i)_i
\]
are contractions for both norms
\[
(u_i)_i\mapsto \Tr\left(\sum_i u_i^*u_i\right)^{1/2},
\qquad
(u_i)_i\mapsto \Tr\left(\sum_i u_i u_i^*\right)^{1/2}.
\]
Therefore any decomposition \(x_i=y_i+z_i\) can be replaced by the
diagonal decomposition \(x_i=D y_i+D z_i\) without increasing the
objective.  Once all matrices are diagonal, the objective is
\[
\sum_{k=1}^n \left(\|b(k)\|_2+\|c(k)\|_2\right),
\qquad b(k)+c(k)=a(k),
\]
where \(a(k)=(a_i(k))_i\).  The pointwise infimum is exactly
\(\|a(k)\|_2\), proving the formula.
\end{proof}

\begin{theorem}[Sharp abelian subcase]
\label{thm:abelian}
If \(x_1,\ldots,x_d\in M_n(\C)\) are commuting normal matrices, then
\[
\E_\eps\left\|\sum_{i=1}^d \eps_i x_i\right\|_{S_1^n}
\geq
\frac1{\sqrt2}N(x_1,\ldots,x_d).
\]
The constant \(1/\sqrt2\) is best possible even in this abelian class.
\end{theorem}

\begin{proof}
After unitary conjugacy the \(x_i\)'s are diagonal,
\(x_i=\diag(a_i(1),\ldots,a_i(n))\).  By Lemma \ref{lem:diag},
\[
N(x_1,\ldots,x_d)=\sum_{k=1}^n \|a(k)\|_2.
\]
On the other hand,
\[
\E_\eps\left\|\sum_i \eps_i x_i\right\|_1
=
\sum_{k=1}^n \E_\eps\left|\sum_i \eps_i a_i(k)\right|.
\]
Szarek's sharp scalar \(p=1\) Khintchine lower bound gives
\[
\E_\eps\left|\sum_i \eps_i a_i(k)\right|
\geq \frac1{\sqrt2}\|a(k)\|_2
\]
for every \(k\).  Summing over \(k\) proves the asserted inequality.

Sharpness follows from the scalar two-term choice \(n=1\),
\(x_1=x_2=1\), for which \(N(x_1,x_2)=\sqrt2\) and
\(\E|\eps_1+\eps_2|=1\).
\end{proof}

\begin{theorem}[Sharp arbitrary two-term subcase]
\label{thm:twoterm}
For all \(x_1,x_2\in M_n(\C)\),
\[
\E_{\eps_1,\eps_2}\left\|\eps_1x_1+\eps_2x_2\right\|_{S_1^n}
\geq
\frac1{\sqrt2}N(x_1,x_2).
\]
The constant \(1/\sqrt2\) is best possible.
\end{theorem}

\begin{proof}
Since the trace norm is invariant under multiplication by \(-1\),
\[
\E_{\eps_1,\eps_2}\|\eps_1x_1+\eps_2x_2\|_1
=\frac12\bigl(\|x_1+x_2\|_1+\|x_1-x_2\|_1\bigr).
\]
Let \(R=[x_1\ x_2]\) be the row block operator.  Then
\[
\|R\|_{S_1}=\Tr(x_1x_1^*+x_2x_2^*)^{1/2}.
\]
Right multiplication by the \(2\times2\) Hadamard unitary gives
\[
\|R\|_{S_1}
=
\frac1{\sqrt2}\|[x_1+x_2\quad x_1-x_2]\|_{S_1}
\leq
\frac1{\sqrt2}\bigl(\|x_1+x_2\|_1+\|x_1-x_2\|_1\bigr).
\]
In the definition of \(N\), taking \(y_1=y_2=0\) and \(z_i=x_i\)
gives
\[
N(x_1,x_2)\leq \Tr(x_1x_1^*+x_2x_2^*)^{1/2}=\|R\|_{S_1}.
\]
Combining the last three displays yields
\[
N(x_1,x_2)
\leq
\sqrt2\,
\E_{\eps_1,\eps_2}\|\eps_1x_1+\eps_2x_2\|_1,
\]
which is the claim.  The scalar example \(x_1=x_2=1\) again gives
equality.
\end{proof}

\section*{Verification notes}

The proof is deterministic.  The included script
\texttt{code/verify\_subcases.py} numerically checks the diagonal norm
formula and the two-term row-block estimate on random small matrices.
These checks are only sanity tests, not part of the proof.

\section*{Limitations and novelty check}

This is not a full solution of the Haagerup--Musat Rademacher sharpness
problem.  It leaves open the case of three or more genuinely
noncommuting coefficients.  A local corpus search found Pisier's 2011
survey text still recording the Rademacher best constants as apparently
unknown \cite{PisierSurvey}.  A 2025 preprint by Chuah--Liu--Mei also
states that the exact Rademacher operator-valued value is still unknown
while proving related \(Z_2\)-set criteria \cite{CLM}.  No exact
literature packet for arXiv:math/0611160 was present in the run
registry at promotion time.

\section*{Human-review recommendation}

Check Lemma \ref{lem:diag}, especially the use of the diagonal
conditional expectation as a contraction on the two \(S_1(\ell_2)\)
row/column norms.  The two-term argument is a short block-matrix
calculation and should be straightforward to verify.

\begin{thebibliography}{9}
\bibitem{HM}
U. Haagerup and M. Musat,
\emph{On the best constants in noncommutative Khintchine-type inequalities},
J. Funct. Anal. 250 (2007), 588--624; arXiv:math/0611160.

\bibitem{Szarek}
S. J. Szarek,
\emph{On the best constants in the Khinchin inequality},
Studia Math. 58 (1976), 197--208.

\bibitem{PisierSurvey}
G. Pisier,
\emph{Grothendieck's theorem, past and present},
Bull. Amer. Math. Soc. 49 (2012), 237--323; arXiv:1101.4195.

\bibitem{CLM}
C. Y. Chuah, Z.-C. Liu, and T. Mei,
\emph{Noncommutative Khintchine inequalities and \(Z_2\) property},
arXiv:2503.03187.
\end{thebibliography}

\end{document}
